\section*{Capítulo \ref{ch:b2:powerseries} --- Series de potencias}

\begin{solution}{exo:b2:powerseries:1}
$\sum \frac{n^2}{2^n}z^n$: cociente $\frac{(n+1)^2}{2^{n+1}}\cdot
\frac{2^n}{n^2} \to \frac12$: $R = 2$.

$\sum \frac{z^n}{\binom{2n}{n}}$: $\binom{2n}{n} \sim
\frac{4^n}{\sqrt{\pi n}}$
(\cref{ex:b2:comparison:centralbinomial}), de modo que
$\abs{a_n}^{-1} \approx 4^n$ salvo factores polinómicos: $R = 4$
(criterio del cociente: $\frac{\binom{2n}{n}}{\binom{2n+2}{n+1}} =
\frac{(n+1)^2}{(2n+1)(2n+2)} \to \frac14$).

$\sum z^{n!}$: coeficientes $a_k = 1$ si $k = n!$ y $0$ en otro
caso. Para $\abs z < 1$, $\sum \abs z^{n!}$ converge (dominada por
una geométrica); y para $\abs z \geq 1$ los términos no tienden a
$0$: $R = 1$.

$\sum (2 + (-1)^n)^n z^n$: coeficientes $3^n$ ($n$ par) y $1$ ($n$
impar). La acotación de $a_nr^n$ exige $3r \leq 1$; y todo $r <
\frac13$ funciona: $R = \frac13$.
\end{solution}

\begin{solution}{exo:b2:powerseries:2}
Las tres tienen radio $1$ (criterio del cociente). En $z = 1$: $\sum
1$ diverge; $\sum\frac1n$ diverge; $\sum\frac{1}{n^2}$ converge. En
$z = -1$: $\sum(-1)^n$ diverge; $\sum\frac{(-1)^n}{n}$ converge
(alternada); $\sum\frac{(-1)^n}{n^2}$ converge (\omterm{def:b2:series:def}{absolutamente}). El
comportamiento en la frontera es invisible para el radio.
\end{solution}

\begin{solution}{exo:b2:powerseries:3}
De $\frac{1}{1-x} = \sum x^n$, derivando y multiplicando por $x$
(\cref{thm:b2:powerseries:calculus}):
\[
\sum n x^n = \frac{x}{(1-x)^2} .
\]
Derivando una vez más y multiplicando de nuevo por $x$:
\[
\sum n^2 x^n = x\,\frac{\dd}{\dd x}\Bigl(\frac{x}{(1-x)^2}\Bigr)
= \frac{x(1 + x)}{(1-x)^3} .
\]
Tercera suma: es la parte impar de $-\ln(1 - x)$:
\[
\sum_{n\geq0} \frac{x^{2n+1}}{2n+1}
= \frac{-\ln(1-x) + \ln(1+x)}{2}
= \frac12 \ln\frac{1+x}{1-x}
= \operatorname{artanh} x .
\]
\end{solution}

\begin{solution}{exo:b2:powerseries:4}
Fracciones simples: $\frac{1}{(1-x)(2-x)} = \frac{1}{1-x} -
\frac{1}{2 - x} = \sum x^n - \frac12\sum \bigl(\frac x2\bigr)^n$:
\[
\frac{1}{(1-x)(2-x)} = \sum_{n\geq0} \Bigl(1 -
\frac{1}{2^{n+1}}\Bigr)x^n,
\qquad R = 1 .
\]

$\ln(1 + x + x^2) = \ln\frac{1 - x^3}{1 - x} = \ln(1 - x^3) - \ln(1
- x) = \sum_{n\geq1}\frac{x^n}{n} -
\sum_{m\geq1}\frac{x^{3m}}{m}$: el coeficiente de $x^n$ es
$\frac1n$ si $3 \nmid n$, y $\frac1n - \frac{3}{n} = -\frac2n$ si $3
\mid n$. Radio $1$ (la obstrucción más próxima: la serie de $\ln(1 -
x^3)$).
\end{solution}

\begin{solution}{exo:b2:powerseries:5}
\omterm{thm:b2:series:fubini}{Producto de Cauchy} de $\sum_{m \geq 0} x^m$ (coeficientes $1$) y
$\sum_{k\geq1} \frac{x^k}{k}$ (coeficientes $\frac1k$, $k \geq 1$),
ambas \omterm{def:b2:series:def}{absolutamente} convergentes para $\abs x < 1$: el
coeficiente de $x^n$ en el producto es $\sum_{k=1}^{n} \frac1k \cdot
1 = H_n$. Por tanto
\[
\Bigl(\sum x^m\Bigr)\Bigl(\sum \frac{x^k}{k}\Bigr)
= \frac{1}{1-x}\cdot\bigl(-\ln(1-x)\bigr)
= \sum_{n\geq1} H_n x^n .
\]
\end{solution}

\begin{solution}{exo:b2:powerseries:6}
Multipliquemos la recurrencia por $x^{n+1}$ y sumemos (para $\abs x
< \frac12$):
\[
U(x) - 1 = 2x\,U(x) + \sum_{n\geq0} n\,x^{n+1}
= 2x\,U(x) + \frac{x^2}{(1-x)^2} ,
\]
usando el~\cref{exo:b2:powerseries:3}. Por tanto
\[
U(x) = \frac{1}{1 - 2x}\Bigl(1 + \frac{x^2}{(1-x)^2}\Bigr)
= \frac{1 - 2x + 2x^2}{(1-2x)(1-x)^2} .
\]
Fracciones simples (tapando en $x = \frac12$ sale el coeficiente
$2$; en el polo doble $x = 1$, el coeficiente $-1$; y el coeficiente
central se anula evaluando en $x = 0$):
\[
U(x) = \frac{2}{1-2x} - \frac{1}{(1 - x)^2} .
\]
Desarrollando ambos:
\[
u_n = 2\cdot 2^n - (n + 1) = 2^{n+1} - n - 1 .
\]
(Comprobación: $u_0 = 1$, $u_1 = 2u_0 + 0 = 2 = 4 - 2$.)
\end{solution}

\begin{solution}{exo:b2:powerseries:7}
Inducción: $f'(x) = \frac{2}{x^3}\eu^{-1/x^2}$ y, si $f^{(n)}(x) =
P_n(\tfrac1x)\eu^{-1/x^2}$, entonces
\[
f^{(n+1)}(x) =
\Bigl(-\frac{1}{x^2}\,P_n'\Bigl(\frac1x\Bigr) +
\frac{2}{x^3}\,P_n\Bigl(\frac1x\Bigr)\Bigr)\eu^{-1/x^2} :
\]
de nuevo de la forma indicada. En $0$: los cocientes incrementales
$\frac{f^{(n)}(h)}{h} = \frac1h P_n(\frac1h)\eu^{-1/h^2} \to 0$
cuando $h \to 0$, ya que $Q(u)\,\eu^{-u^2} \to 0$ cuando $u \to
\pm\infty$ para todo polinomio $Q$ (la exponencial gana a las
potencias): por inducción, todas las $f^{(n)}(0)$ existen y se
anulan, y cada $f^{(n)}$ es \omterm{def:b2:metric:continuity}{continua} en $0$ por el mismo límite.
Así pues, $f \in C^\infty$ con serie de Taylor nula en $0$; esa
serie suma $0 \neq f$: no es \omterm{def:b2:powerseries:analytic}{analítica} en $0$.
\end{solution}

\begin{solution}{exo:b2:powerseries:8}
Serie binomial: $\sqrt{1-4x} = \sum_{k\geq0}
\binom{1/2}{k}(-4x)^k$. Para $k = n + 1 \geq 1$:
\begin{align*}
\binom{1/2}{n+1}(-4)^{n+1}
&= \frac{\frac12\bigl(\frac12 - 1\bigr)\cdots\bigl(\frac12 -
n\bigr)}{(n+1)!}\,(-4)^{n+1}\\
&= \frac{(-1)^n\,1\cdot3\cdots(2n-1)}{2^{n+1}(n+1)!}\,(-4)^{n+1}
= -\frac{2}{n+1}\cdot\frac{(2n)!}{n!\,n!} ,
\end{align*}
usando $1\cdot3\cdots(2n-1) = \frac{(2n)!}{2^n n!}$. Por tanto
\[
C(x) = \frac{1 - \sqrt{1-4x}}{2x}
= \frac{1}{2x}\sum_{n\geq0}\frac{2}{n+1}\binom{2n}{n}x^{n+1}
= \sum_{n\geq0} \frac{1}{n+1}\binom{2n}{n}\,x^n :
\]
$C_n = \frac{1}{n+1}\binom{2n}{n}$. Radio: $\frac14$ (la serie
binomial en $4x$). Asintótica vía el~\cref{ex:b2:comparison:centralbinomial}:
\[
C_n \sim \frac{4^n}{\sqrt{\pi}\; n^{3/2}} .
\]
\end{solution}

\begin{solution}{exo:b2:powerseries:9}
Con $A_n = \sum_{k \leq n} a_k \to A$, la \omterm{thm:b2:series:abel}{sumación de Abel} da, para
$0 \leq x < 1$,
\[
\sum_{n=0}^{\infty} a_n x^n = (1 - x)\sum_{n=0}^{\infty} A_n x^n
\]
(ambos miembros convergen, pues $(A_n)$ está acotada; y la identidad
se sigue de $a_n = A_n - A_{n-1}$ y de reindexar). Como $(1 -
x)\sum x^n = 1$:
\[
\sum_n a_nx^n - A = (1-x)\sum_{n} (A_n - A)x^n .
\]
Dado $\varepsilon$, elíjase $N$ con $\abs{A_n - A} \leq
\varepsilon$ para $n > N$; entonces
\[
\Bigl|\sum a_nx^n - A\Bigr|
\leq (1-x)\sum_{n \leq N}\abs{A_n - A} + \varepsilon(1 -
x)\sum_{n > N}x^n
\leq (1-x)\,C_N + \varepsilon ,
\]
y haciendo $x \to 1^-$: el límite superior es $\leq \varepsilon$
para todo $\varepsilon$. Por tanto, el límite radial vale $A$.

Aplicaciones: $\sum \frac{(-1)^{n-1}}{n}$ converge (alternada) y,
para $x < 1$, su serie de potencias suma $\ln(1 + x)$: por Abel, la
suma es $\ln 2$. Análogamente, $\sum\frac{(-1)^n}{2n+1}x^{2n+1} =
\arctan x$ da $\frac\pi4$ en $x = 1$: las demostraciones integrales
del primer año, ahora estructurales.
\end{solution}

\begin{solution}{exo:b2:powerseries:10}
Tanto $\sum\frac{x^n}{n}$ como $\sum\frac{x^n}{n+1}$ tienen radio
$1$, y $\frac{1}{n(n+1)} = \frac1n - \frac1{n+1}$, de modo que para
$0 < \abs x < 1$:
\[
\sum_{n\geq1}\frac{x^n}{n(n+1)}
= -\ln(1-x) - \frac1x\sum_{n\geq1}\frac{x^{n+1}}{n+1}
= -\ln(1-x) - \frac{-\ln(1-x) - x}{x}
= 1 + \frac{1-x}{x}\ln(1-x) .
\]
Radio $1$; y $\norm{x^n/(n(n+1))}_{\infty,\intcc{-1}1} =
\frac{1}{n(n+1)}$ es \omterm{def:b2:series:summable}{sumable}: hay \omterm{def:b2:funcseq:series}{convergencia normal} sobre
$\intcc{-1}{1}$, luego la suma es allí \omterm{def:b2:metric:continuity}{continua}. Cuando $x \to
1^-$, $(1-x)\ln(1-x) \to 0$ y la forma cerrada tiende a $1$, de
acuerdo con el valor telescópico $\sum\frac{1}{n(n+1)} = \lim_N\bigl(1
- \frac{1}{N+1}\bigr) = 1$ en $x = 1$.
\end{solution}

\begin{solution}{exo:b2:powerseries:11}
Clasificando las $n!$ permutaciones por su conjunto de puntos fijos:
elegir los $k$ puntos fijos ($\binom nk$ maneras) y desarreglar los
otros $n - k$ objetos da $n! = \sum_{k=0}^n\binom nk D_{n-k}$.
Dividiendo por $n!$:
\[
1 = \sum_{k=0}^{n}\frac{1}{k!}\cdot\frac{D_{n-k}}{(n-k)!} ,
\]
que dice exactamente que el \omterm{thm:b2:series:fubini}{producto de Cauchy} de $\eu^x =
\sum\frac{x^k}{k!}$ y $D(x) = \sum D_n\frac{x^n}{n!}$ es $\sum x^n =
\frac{1}{1-x}$. Ambos factores convergen \omterm{def:b2:series:def}{absolutamente} para $\abs x
< 1$ ($D_n \leq n!$, de modo que $D$ está dominada por la serie
geométrica): la identidad del producto es legítima
(\cref{prop:b2:powerseries:operations}) y
\[
D(x) = \frac{\eu^{-x}}{1-x} .
\]
\omterm{thm:b2:series:fubini}{Producto de Cauchy} de $\eu^{-x} = \sum\frac{(-1)^kx^k}{k!}$ y
$\sum x^m$: el coeficiente de $x^n$ es
$\sum_{k=0}^{n}\frac{(-1)^k}{k!}$ y, por la unicidad de los
coeficientes de una serie de potencias
(\cref{thm:b2:powerseries:calculus}):
\[
\frac{D_n}{n!} = \sum_{k=0}^{n}\frac{(-1)^k}{k!}
\xrightarrow[n\to\infty]{} \eu^{-1} :
\]
alrededor del $37\,\%$ de todas las permutaciones son desarreglos,
sea cual sea $n$.
\end{solution}

\begin{solution}{exo:b2:powerseries:12}
Serie binomial con $\alpha = -\frac12$ en $-4x$:
\[
\binom{-1/2}{n}(-4)^n
= \frac{\bigl(-\frac12\bigr)\bigl(-\frac32\bigr)\cdots
\bigl(-\frac{2n-1}2\bigr)}{n!}(-4)^n
= \frac{1\cdot3\cdots(2n-1)}{2^n\,n!}\,4^n
= \frac{(2n)!}{2^n n!}\cdot\frac{2^n}{n!}
= \binom{2n}{n},
\]
usando $1\cdot3\cdots(2n-1) = \frac{(2n)!}{2^nn!}$. Por tanto,
$(1-4x)^{-1/2} = \sum\binom{2n}nx^n$ para $\abs{4x} < 1$. Elevando
al cuadrado (\omterm{thm:b2:series:fubini}{producto de Cauchy}, legítimo por la convergencia
absoluta) y comparando con $\frac{1}{1-4x} = \sum 4^nx^n$: el
coeficiente de $x^n$ en el cuadrado es
$\sum_{k=0}^n\binom{2k}k\binom{2n-2k}{n-k}$, y la unicidad de los
coeficientes da
\[
\sum_{k=0}^{n}\binom{2k}{k}\binom{2n-2k}{n-k} = 4^n .
\]
\end{solution}

\begin{solution}{pb:b2:powerseries:1}
\textbf{1.} Con $a_n = r_n - r_{n+1}$, la sumación por partes da
\[
\sum_{n=N}^{M} a_nx^n
= r_Nx^N + \sum_{n=N+1}^{M} r_n\bigl(x^n - x^{n-1}\bigr)
- r_{M+1}x^M .
\]
Para $0 \leq x \leq 1$, los incrementos $x^{n-1} - x^n$ son no
negativos y telescopan a $x^N - x^M$; con $s = \sup_{n\geq
N}\abs{r_n}$:
\[
\Bigl|\sum_{n=N}^{M}a_nx^n\Bigr|
\leq s\bigl(x^N + (x^N - x^M) + x^M\bigr) = 2s\,x^N \leq 2s .
\]

\textbf{2.} Como $r_n \to 0$, $\sup_{n\geq N}\abs{r_n} \to 0$: la
pregunta 1 es exactamente el criterio de Cauchy uniforme sobre
$\intcc{0}{1}$, de modo que $\sum a_nx^n$ converge allí
\omterm{def:b2:funcseq:def}{uniformemente} y su suma es \omterm{def:b2:metric:continuity}{continua}
(\cref{thm:b2:funcseq:seriestransfer}). Siendo el valor en $1$
igual a $\sum a_n$, la \omterm{def:b2:metric:continuity}{continuidad} en $1$ es el límite radial del~\cref{exo:b2:powerseries:9}.

\textbf{3.} Para $\abs x < 1$, las tres series de potencias
convergen \omterm{def:b2:series:def}{absolutamente} y $\bigl(\sum a_nx^n\bigr)\bigl(\sum
b_nx^n\bigr) = \sum c_nx^n$
(\cref{prop:b2:powerseries:operations}). Por la pregunta 2, cada
factor y el miembro del producto son \omterm{def:b2:metric:continuity}{continuos} sobre
$\intcc{0}{1}$ (sus series de coeficientes convergen por
hipótesis); haciendo $x \to 1^-$ en la identidad: $AB = C$.

\textbf{4.} Aquí
\[
\abs{c_n} = \sum_{k=0}^{n}
\frac{1}{\sqrt{(k+1)(n-k+1)}}
\geq \sum_{k=0}^{n}\frac{2}{n+2}
= \frac{2(n+1)}{n+2} \geq 1,
\]
por la desigualdad entre medias: $\sqrt{(k+1)(n-k+1)} \leq
\frac{(k+1) + (n-k+1)}{2} = \frac{n+2}{2}$. El término general de
$\sum c_n$ no tiende a $0$: el \omterm{thm:b2:series:fubini}{producto de Cauchy} diverge, pese a
que ambos factores convergen (series alternadas).

\textbf{5.} El~\cref{exo:b2:powerseries:5} da $\frac{-\ln(1-x)}{1-x}
= \sum H_nx^n$ ($\abs x < 1$). Primitivas término a término
(\cref{thm:b2:powerseries:calculus}~(2)), anulándose ambos miembros
en $0$:
\[
\frac{\bigl(\ln(1-x)\bigr)^2}{2}
= \sum_{n\geq1}\frac{H_n}{n+1}\,x^{n+1} .
\]
Decrecimiento: $(n+2)H_n \geq (n+1)H_{n+1}$ equivale a $H_n \geq
1$, cierto para $n \geq 1$; y $\frac{H_n}{n+1} \sim \frac{\ln n}{n}
\to 0$: en $x = -1$ la serie converge por el criterio de las
alternadas. Sustituyendo $x \mapsto -x$ y aplicando la pregunta 2 en
$x = 1$:
\[
\frac{(\ln 2)^2}{2}
= \sum_{n\geq1}\frac{H_n}{n+1}(-1)^{n+1},
\]
el valor anunciado.

\textbf{6.} $f(x) = \sum(-1)^nx^n = \frac{1}{1+x} \to \frac12$
cuando $x \to 1^-$: \omterm{def:b2:series:summable}{Abel-sumable} a $\frac12$. Pero las sumas
parciales son $1, 0, 1, 0, \dots$: divergente.

\textbf{7.} Dado $\varepsilon > 0$, elíjase $m$ con $\abs{u_n} \leq
\varepsilon$ para $n > m$; para $N \geq m$:
\[
\Bigl|\frac{u_1 + \dots + u_N}{N}\Bigr|
\leq \frac{\abs{u_1} + \dots + \abs{u_m}}{N} +
\varepsilon\,\frac{N - m}{N}
\leq \frac{C_m}{N} + \varepsilon,
\]
luego el límite superior es $\leq \varepsilon$ para todo
$\varepsilon$: las medias tienden a $0$.

\textbf{8.} Para $0 \leq x \leq 1$: $1 - x^n = (1-x)(1 + x + \dots +
x^{n-1}) \leq n(1-x)$, luego
\[
\Bigl|\sum_{n=0}^N a_n(1 - x_N^n)\Bigr|
\leq (1 - x_N)\sum_{n=1}^N n\abs{a_n}
= \frac1N\sum_{n=1}^{N}n\abs{a_n} .
\]
Para $n > N$: $\abs{a_n} = \frac{n\abs{a_n}}{n} \leq
\frac{1}{N}\sup_{m>N}m\abs{a_m}$, y $\sum_{n>N}x_N^n \leq \frac{1}{1
- x_N} = N$:
\[
\Bigl|\sum_{n>N}a_nx_N^n\Bigr|
\leq \frac{\sup_{m>N}m\abs{a_m}}{N}\cdot N
= \sup_{m>N}\,m\abs{a_m} .
\]

\textbf{9.} Descompongamos
\[
A_N - L = \sum_{n=0}^{N}a_n\bigl(1 - x_N^n\bigr)
- \sum_{n>N}a_nx_N^n + \bigl(f(x_N) - L\bigr) .
\]
El primer término tiende a $0$ por la pregunta 7 (las medias de
$n\abs{a_n} \to 0$), el segundo por la pregunta 8 (el supremo tiende
a $0$) y el tercero porque $x_N \to 1^-$ y $f(x) \to L$. Por tanto
$A_N \to L$: el teorema de Tauber.

\textbf{10.} Para $x \in \intco{0}{1}$ y todo $N$: $\sum_{n\leq
N}a_nx^n \leq f(x) \leq M$ (términos no negativos). Haciendo $x \to
1^-$ en la suma finita: $\sum_{n\leq N}a_n \leq M$. Las sumas
parciales son crecientes y acotadas: $\sum a_n$ converge, y entonces
la pregunta 2 da $\lim_{x\to1^-}f(x) = \sum a_n$.

\textbf{11.} $\sigma_N - L$ es la media de los $N$ números $A_n - L$
($0 \leq n < N$), que tienden a $0$: pregunta 7.

\textbf{12.} $A_n = 1$ para $n$ par y $0$ para $n$ impar: $A_0 +
\dots + A_{N-1} = \lceil N/2\rceil$, de modo que $\sigma_N =
\frac{\lceil N/2\rceil}{N} \to \frac12$, el valor de Abel de la
pregunta 6.

\textbf{13.} Bajo $\sigma_N \to L$ se tiene $S_n = O(n)$, de donde
$A_n = S_n - S_{n-1} = O(n)$ y $a_n = O(n)$: todas las series de más
abajo tienen radio $\geq 1$. Para $\abs x < 1$, de $a_n = A_n -
A_{n-1}$ y $A_nx^n \to 0$:
\[
(1-x)\sum_n A_nx^n = \sum_n A_nx^n - \sum_n A_nx^{n+1}
= \sum_n a_nx^n = f(x),
\]
e idénticamente $(1-x)\sum S_nx^n = \sum A_nx^n$, luego $f(x) =
(1-x)^2\sum_n S_nx^n = (1-x)^2\sum_n(n+1)\sigma_{n+1}x^n$. Por
último, $\sum(n+1)x^n = \frac{1}{(1-x)^2}$
(\cref{exo:b2:powerseries:3}), que es la segunda identidad.

\textbf{14.} Restando $L$ veces la segunda identidad de la primera:
\[
f(x) - L = (1-x)^2\sum_{n\geq0}(n+1)
\bigl(\sigma_{n+1} - L\bigr)x^n .
\]
Dado $\varepsilon$, elíjase $N$ con $\abs{\sigma_{n+1} - L} \leq
\varepsilon$ para $n \geq N$; entonces
\[
\abs{f(x) - L} \leq (1-x)^2 C_N +
\varepsilon(1-x)^2\sum_{n}(n+1)x^n
= (1-x)^2C_N + \varepsilon ,
\]
y haciendo $x \to 1^-$: el límite superior es $\leq \varepsilon$.
Por tanto $f(x) \to L$: el teorema de Frobenius.

\textbf{15.} $f(x) = \sum(-1)^n(n+1)x^n = \frac{1}{(1+x)^2}$
(derívese la serie geométrica en $-x$): valor de Abel $\frac14$.
Sumas parciales: $A_{2k} = k+1$ y $A_{2k+1} = -(k+1)$ (inducción
inmediata). Entonces $S_{2m-1} = 0$ (los pares consecutivos se
cancelan) y $S_{2m} = m + 1$, de modo que
\[
\sigma_{2m} = \frac{S_{2m-1}}{2m} = 0,
\qquad
\sigma_{2m+1} = \frac{m+1}{2m+1} \to \frac12 :
\]
$(\sigma_N)$ tiene dos valores de acumulación distintos: no es
\omterm{def:b2:series:summable}{Cesàro-sumable}. Con las preguntas 6 y 11--14, la jerarquía
convergente $\Rightarrow$ Cesàro $\Rightarrow$ Abel es estricta en
ambas flechas.

\textbf{16.} La serie primitiva $F(x) = \sum\frac{a_n}{n+1}x^{n+1}$
tiene el mismo radio y $F' = f$ sobre $\intco{0}{1}$
(\cref{thm:b2:powerseries:calculus}); y como $\sum\frac{a_n}{n+1}$
converge, la parte I (pregunta 2) hace $F$ \omterm{def:b2:metric:continuity}{continua} sobre
$\intcc{0}{1}$. Como $\int_0^x f = F(x)$ (derivadas iguales, valor
$0$ en $0$),
\[
\int_0^x f \xrightarrow[x\to1^-]{} F(1)
= \sum_{n\geq0}\frac{a_n}{n+1} :
\]
la \omterm{def:b2:integration:improper}{integral impropia} existe con el valor indicado.

\textbf{17.} $\frac{\ln(1+x)}{x} =
\sum_{n\geq1}\frac{(-1)^{n-1}}{n}x^{n-1}$ (radio $1$; \omterm{def:b2:metric:continuity}{continua} en
$0$). La serie de los $\frac{a_m}{m+1}$ es
$\sum_{n\geq1}\frac{(-1)^{n-1}}{n^2}$, \omterm{def:b2:series:def}{absolutamente} convergente:
la pregunta 16 da $\int_0^1\frac{\ln(1+x)}{x}\dd x = \eta$. En la
serie \omterm{def:b2:series:def}{absolutamente} convergente $\sum\frac1{n^2}$, reagrupemos pares
e impares:
\[
\eta = \sum_{\text{impares}}\frac1{n^2} -
\sum_{\text{pares}}\frac1{n^2}
= \sum_{n}\frac1{n^2} - 2\sum_{k}\frac1{(2k)^2}
= \Bigl(1 - \frac12\Bigr)\sum_n\frac1{n^2}
= \frac12\sum_{n\geq1}\frac1{n^2} .
\]

\textbf{18.} Leibniz: $\frac{1}{3n+1}\downarrow0$, luego la serie
converge. Sobre $\intco{0}{1}$, $\sum(-1)^nx^{3n} =
\frac{1}{1+x^3}$, y $\sum\frac{(-1)^n}{3n+1}$ converge: la pregunta
16 da $\sum\frac{(-1)^n}{3n+1} = \int_0^1\frac{\dd x}{1+x^3}$.
Fracciones simples (compruébese: $\frac13(x^2-x+1) +
\frac{2-x}{3}(1+x) = 1$):
\[
\int_0^1\frac{\dd x}{1+x^3}
= \frac13\ln2 + \frac13\int_0^1\frac{2-x}{x^2-x+1}\dd x .
\]
Escribiendo $2 - x = -\frac12(2x-1) + \frac32$: la parte en
$\ln(x^2-x+1)$ se anula en ambos extremos, y
\[
\frac32\int_0^1\frac{\dd x}{(x-\frac12)^2 + \frac34}
= \frac32\cdot\frac{2}{\sqrt3}
\Bigl[\arctan\frac{2x-1}{\sqrt3}\Bigr]_0^1
= \sqrt3\cdot\frac{\pi}{3} = \frac{\pi}{\sqrt3} .
\]
Total: $\frac13\bigl(\ln2 + \frac{\pi}{\sqrt3}\bigr)$.

\textbf{19.} Sustituyendo $t = x^2$ en la serie del~\cref{exo:b2:powerseries:12} e integrando término a término (la
primitiva de $(1-x^2)^{-1/2}$ que se anula en $0$ es $\arcsin$):
\[
\arcsin x = \sum_{n\geq0}
\frac{\binom{2n}n}{4^n(2n+1)}x^{2n+1}
\qquad(\abs x < 1) .
\]
Los coeficientes son $\sim \frac{1}{2\sqrt\pi\,n^{3/2}}$
(\cref{ex:b2:comparison:centralbinomial}), \omterm{def:b2:series:summable}{sumables}: la serie
converge \emph{\omterm{def:b2:funcseq:series}{normalmente}} sobre $\intcc{-1}{1}$, su suma es allí
\omterm{def:b2:metric:continuity}{continua} y coincide con la función \omterm{def:b2:metric:continuity}{continua} $\arcsin$ sobre
$\intoo{-1}{1}$ y, por tanto, también en $x = 1$:
\[
\sum_{n\geq0}\frac{\binom{2n}n}{4^n(2n+1)}
= \arcsin 1 = \frac\pi2 .
\]

\textbf{20.} $C_n4^{-n} \sim \frac{1}{\sqrt\pi\,n^{3/2}}$
(\cref{exo:b2:powerseries:8}): hay \omterm{def:b2:funcseq:series}{convergencia normal} de $\sum
C_nx^n$ sobre $\intcc{0}{\frac14}$, luego su suma es allí
\omterm{def:b2:metric:continuity}{continua}; sobre $\intoo{0}{\frac14}$ vale
$\frac{1-\sqrt{1-4x}}{2x}$ (\cref{ex:b2:powerseries:catalan}), cuyo
límite en $\frac14^-$ es $\frac{1-0}{1/2} = 2$. Por tanto,
$\sum_{n\geq0} C_n4^{-n} = 2$.

\textbf{21.} $f(x) = \sum_{n\geq m}a_nx^n = x^m g(x)$ con $g(x) =
\sum_{k\geq0}a_{m+k}x^k$; si $(a_nr^n)$ está acotada, también lo
está $(a_{m+k}r^k)$ (divídase por $r^m$): $g$ tiene radio $\geq R$.
Y $g$ es \omterm{def:b2:metric:continuity}{continua} con $g(0) = a_m \neq 0$, de modo que $g \neq 0$
sobre algún $\intcc{-\delta}{\delta}$ y $f(x) = x^mg(x) \neq 0$ para
$0 < \abs x \leq \delta$.

\textbf{22.} $d = f - h$ es la suma de una serie de potencias cerca
de $0$ que se anula en los puntos no nulos $x_k \to 0$. Si algún
coeficiente de $d$ fuera no nulo, la pregunta 21 daría un entorno
perforado de $0$ libre de ceros de $d$, en contra de $d(x_k) = 0$.
Así pues, todos los coeficientes de $d$ se anulan: $f$ y $h$ tienen
los mismos coeficientes y coinciden cerca de $0$.

\textbf{23.} La función $h(x) = \frac{1}{1+x^2} = \sum(-1)^nx^{2n}$
(radio $1$) cumple $h(\frac1k) = \frac{1}{1 + 1/k^2} =
\frac{k^2}{k^2+1}$. Toda $f$ \omterm{def:b2:powerseries:analytic}{analítica} con los mismos valores
coincide con $h$ en los puntos $\frac1k \to 0$: por el teorema de
identidad (pregunta 22), $f = \frac{1}{1+x^2}$ cerca de $0$; la
solución única.

\textbf{24.} Sea $Z$ el conjunto de los puntos de $I$ que tienen un
entorno sobre el cual $f$ se anula idénticamente: es \omterm{def:b2:metric:topology}{abierto} por
definición y no vacío (el subintervalo). Cerrado en $I$: si $y \in
I$ es límite de puntos de $Z$, entonces $y$ es punto de acumulación
de ceros de $f$; desarrollando $f$ en serie de potencias en $y$ (por
analiticidad) y aplicando las preguntas 21--22 recentradas en $y$,
todos los coeficientes en $y$ se anulan, luego $f \equiv 0$ cerca de
$y$: $y \in Z$. Un intervalo es \omterm{def:b2:metric:connected}{conexo}, así que $Z = I$: $f \equiv
0$ sobre $I$. En particular, una \omterm{def:b2:powerseries:analytic}{función analítica} sobre $\R$ que se
anula fuera de un \omterm{def:b2:metric:compact}{compacto} se anula sobre un intervalo y, por tanto,
en todas partes: no hay mesetas \omterm{def:b2:powerseries:analytic}{analíticas} no nulas. El mundo
$C^\infty$ es distinto: pegando la función plana del~\cref{exo:b2:powerseries:7} (por ejemplo, $x \mapsto
\eu^{-1/x^2}\mathbf 1_{x>0}$ y su reflejada) se obtienen mesetas
suaves de soporte \omterm{def:b2:metric:compact}{compacto}.

\textbf{25.} (i) La \omterm{def:b2:funcseq:series}{convergencia normal} vive sobre subdiscos
\omterm{def:b2:metric:compact}{compactos} estrictamente interiores al disco; el teorema de Abel
extiende la \omterm{def:b2:metric:continuity}{continuidad} a un punto de la frontera bajo la única
hipótesis de que la serie de coeficientes converja allí. (ii) El
recíproco vale bajo la condición de Tauber $na_n \to 0$ (pregunta
9), y la \omterm{def:b2:series:summable}{sumabilidad} de Cesàro se sitúa estrictamente entre la
convergencia y la \omterm{def:b2:series:summable}{sumabilidad} de Abel (Frobenius, preguntas
14--15). (iii) Para un amigo: $\sum\frac{(-1)^n}{3n+1} =
\int_0^1\frac{\dd x}{1+x^3}$, integrando la serie geométrica hasta
la frontera y aplicando después fracciones simples. (iv) Las medias
de Cesàro regresan en el capítulo de Fourier como el teorema de
Fejér, donde promediar las sumas parciales repara el fallo de la
\omterm{def:b2:funcseq:def}{convergencia puntual}: misma medicina, paciente nuevo.
\end{solution}
